\section{Background information}
The machine learning market has been developing rapidly since 2020, with the market value being projected to be around 282 billion dollars in size. Because of this rapid development, there is a larger focus on the usage of these models and the harm they can cause.

In 2020, the documentary, "Coded Bias" was released on Netflix. This documentary focuses on the research of MIT media researcher Joy Buolamwini. Her research focuses on the data collection issue plaguing machine learning models' ability to not discriminate. This research started because facial recognition algorithms were found to have significant difficulties in identifying faces of people of color and women. \cite{CodedBias}

The documentary not only covers Buolamwini's research but also the impacts of machine learning models. One notorious example is the group Big Brother Watch in the UK which studies how the biased facial recognition systems trigger law enforcement systems with false positives and cause major stress for individuals. \cite{BigBrotherWatch}

Lastly, the documentary explores how there could be possible improvements for both the algorithms and the structures that produce the models. There is a heavy emphasis on the lack of active regulation and how the law needs to keep up with the technology being developed to prevent further issues from arising. \cite{LegalDefenseFund}

\subsection{Stakeholders}
There are several stakeholders in this scenario with a major power imbalance between the machine learning companies and the general public:

\begin{itemize}
	\item Machine Learning Companies: As they decide what gets made and how it is made
	\item Machine Learning Engineers: As they implement the algorithms and methods of training
	\item Governments: As they regulate and use the algorithms for various purposes
	\item Software companies: As they use the algorithms to enhance their businesses
	\item General public: As they are the ones who are being judged by the algorithm
\end{itemize}

This power imbalance places the general public at a major disadvantage as machine learning algorithms rely on pre-existing data which is not guaranteed to be sanitized of biases. Therefore causing unjust profiling of individuals, rejections in jobs because of nationality/ethnicity, and more without any form of control as the mathematical models are a black box and companies have no incentive to reverse-engineer it.

\section{Salient Issues}
The documentary's main argument is that in computer systems that rely on machine learning are not accounted for in the USA's legal system at a reasonable pace. Therefore allowing companies to neglect the possible biases that the algorithms might have. This can and has been causing violations of human rights as different industries adopt this technology to handle decision-making.\cite{BigBrotherWatch}

Since the original studies and the documentary, the usage of machine learning has increased across all industries \cite{ML-Growth}. This in turn further pushes the need for boundaries in the development, training and usage of said models by major entities like multi-national companies and governments across the world. In this essay, there will be an attempt to apply ethical frameworks to evaluate the following questions:

\begin{itemize}
	\item What moral duties do ML companies have when producing their models?
	\item Who would be acting unethically for not taking responsibility when the models cause damage?
	\item Is it ethically feasible to enforce regulations in the black box that are ML model?
\end{itemize}
